\documentclass[a4paper,10pt]{article}
\usepackage{paper-en}
\usepackage{acts}
\usepackage{hyperref}

%\usepackage[notref,notcite,color]{showkeys}

\def\thetitle{4-string real braid space is CAT(0)}
\def\theauthors{Dmitri Panov and Anton Petrunin}

\hypersetup{colorlinks=true,
citecolor=black,
linkcolor=black,
anchorcolor=black,
filecolor=black,
menucolor=black,
urlcolor=black,
pdftitle={\thetitle},
pdfauthor={\theauthors}
}

%\usepackage[a-2b,mathxmp]{pdfx}[2018/12/22]
%\overfullrule=100mm
%\usepackage[none]{hyphenat}

%\makeatletter\let\orig@section\section\renewcommand{\section}[1]{\orig@section{#1.}}\usepackage{environ}\RenewEnviron{figure}[1][]{}\RenewEnviron{figure*}[1][]{}\RenewEnviron{wrapfigure}[1][]{}\pagestyle{empty}\let\ps@plain\ps@empty\makeatother\def\emph{\textit}\usepackage[none]{hyphenat}

\begin{document}

\title{\thetitle}
\author{\theauthors}
\date{}
\maketitle

\begin{abstract}

\end{abstract}

\section{Introduction}

Let us equip the complex space $\CC^n$ with the standard Euclidean metric.
Let $W$ be $\CC^n$ with removed hyperplanes $z_i=z_j$ for all $i\ne j$; here $z_1,\dots,z_n$ are the standard coordinates.
Denote by $\tilde W$ the universal cover of $W$, equipped with the induced metric.
The completion of $\tilde W$ will be called the \emph{complex braid space}%
\footnote{The terminology comes from the usual interpretation of $W$ as the configuration space of $n$-tuples of distinct points in the complex plane.
A path in $W$ describes a motion of these points without collisions; so their trajectories form a braid with $n$ strings.
The braid group $B_n$ acts isometrically on $\CB^n$, and the quotient space $\CB^n/B_n$ is $\CC^n/S_n$;
here $S_n$ is the symmetry group that acts on $\CC^n$ by permuting coordinates.}
with $n$ strings and denoted by $\CB^n$.

The following conjecture is a special case of a conjecture formulated by Daniel Allcock~\cite[Conjecture 1.4]{allcock}, which is related to an earlier conjecture of Ruth Charney and Michael Davis (see \cite[Conjecture 3]{charney-davis-95}), which in turn is motivated by a conjecture of Arnold, Pham and Thom on complex hyperplane arrangements.
We are responsible for a more general version of Allcock's conjecture \cite[Conjecture 1.1]{panov-petrunin}, which we have been trying to prove for several years.

\begin{thm}{Conjecture}\label{braid-conj}
$\CB^n$ is $\CAT(0)$ for any integer $n\ge 1$.
\end{thm}

The case $n=2$ of \ref{braid-conj} is easy.
The case $n=3$ is already nontrivial.
Two proofs of this statement are given by the authors \cite{panov-petrunin};
one of them follows from the result of Ruth Charney and Michael Davis \cite{charney-davis-93}.
In this note we describe a small step toward the case $n=4$.

Consider the real subspace $\RR^n\subset\CC^n$.
Denote by $\RB^n$ inverse image of $\RR^n$ under the natural projection $\CB^n\to \CC^n$;
we will call it the \emph{real braid space} with $n$ strings.
One may think of the real braid space as a space of braid diagrams.

\begin{thm}{Theorem}\label{thm:main}
$\RB^4$ is $\CAT(0)$.
\end{thm}

By Lemma~\ref{lem:convex-preimage} below, $\RB^n$ is a convex subset of $\CB^n$.
In particular, if $\CB^n$ is $\CAT(0)$, then so is $\RB^n$.
Therefore, the theorems gives a partial case of the conjecture for $n=4$.


\section{Real braid space}

\begin{thm}{Lemma}\label{lem:convex-preimage}
Let $K\subset \CC$ be a closed convex set.
Then the inverse image $\KB^n$ of $K^n\subset \CC^n$ under the natural projection $\CB^n\to \CC^n$ is convex.

In particular, $\RB^n$ is a convex subset of $\CB^n$.
\end{thm}

\parit{Proof.}
The nearest-point projection $r\:\CC\to K$ is short.
Applying $r$ to each strand of a braid gives a short retraction
$\CB^n\to \KB^n$.

Moreover, if a string with endpoints in $\KB^n$ leaves this set, then the retraction decreases its length.
In particular, a minimizing geodesic with endpoints in $\KB^n$ cannot leave $\KB^n$,
and the convexity follows.
Taking $K=\RR$ gives the last assertion.
\qeds

The space $\RB^4$ is a 4-dimensional polyhedral space.
Since all hyperplanes $z_i=z_j$ share the diagonal line $z_1=\ldots=z_n$,
it splits as a Euclidean product space $\RR\times C^3$, where $C^3$ is a 3-dimensional cone over a 2-dimensional spherical polyhedral space $\Sigma$.

The space $\Sigma$ is glued from a countable collection of copies of the spherical triangle $\triangle$ with angles $\tfrac\pi2$, $\tfrac\pi3$, and $\tfrac\pi3$.
Let us describe the gluing.
First denote the sides of $\triangle$ by $a$, $b$ and $c$, so $a$ and $b$ are adjacent to the right angle.
We will keep the notation $a$, $b$ and $c$ for the standard generators of the braid group $B_4$,
so $a\cdot b=b\cdot a$, $a\cdot c\cdot a=c\cdot a\cdot c$ and $b\cdot c\cdot b=c\cdot b\cdot c$.
Take a copy of $\triangle$ for each element of $B_4$ and glue two triangles along one side $a$, $b$, or $c$ if their elements differ by left multiplication by a power of $a$, $b$, or $c$, respectively.

Now let us reduce \ref{thm:main} to the next claim, which will be roved in the following sections.

\begin{thm}{Key claim}\label{clm:sigma}
The space $\Sigma$ does not have a closed geodesic of length $<2\cdot\pi$.
\end{thm}

\parit{Reduction of \ref{thm:main} to \ref{clm:sigma}.}
Note that $\RB^4$ is $\CAT(0)$ if and only if $\Sigma$ is $\CAT(1)$.

The description above implies that $\Sigma$ is locally $\CAT(1)$.
Therefore, by the generalized Cartan--Hadamard theorem, it is sufficient to prove  that
any closed curve in $\Sigma$ of length $<2\cdot\pi$ is null-homotopic in the class of curves of length $<2\cdot\pi$.

Arguing by contradiction, we can find a shortest curve closed curve in $\Sigma$ of length $<2\cdot\pi$ that is not null-homotopic in the class of curves of length $<2\cdot\pi$.
This step is possible since $\Sigma$ admits a cocompact action by $B_4$ and so we can apply the limiting procedure. %???
\qeds



\section{Short trajectories}

\begin{wrapfigure}{o}{36mm}
\centering
\vskip-3mm
\includegraphics{mppics/pic-5}
\vskip-0mm
\end{wrapfigure}

Recall that $\triangle$ denotes the spherical triangle with angles $\tfrac\pi2$, $\tfrac\pi3$, and $\tfrac\pi3$.
Let us denote by $A$, $B$, and $C$ the vertices of $\triangle$ with right angle at $C$,
and let $M$ be the midpoint of the side $AB$.
A closed polygonal line can be encoded by its vertices in the cyclic order; for example, $CM$ denotes the polygonal line that comes from $C$ to $M$, and comes back.

Suppose $\Sigma$ has a closed stable geodesic $\tilde \gamma$ of length less than $2\cdot\pi$;
denote by $\gamma$ its projection to $\triangle$.
Note that $\gamma$ satisfies the following properties:
\begin{itemize}
\item $\gamma$ is a closed polygonal line with all vertices on the boundary of $\triangle$;
\item if a vertex of $\gamma$ lies in an interior of a side of $\triangle$, then $\gamma$ behaves as a billiard trajectory; that is, the angle of incidence is equal to the angle of reflection.
\item if $\gamma$ has a vertex $C$, then it switches the direction at $C$ to the opposite one.
\item If $\gamma$ has a vertex $A$ or $B$, then there is no restriction on the directions of incoming and outgoing edges.
\item Any arc of $\gamma$ that does not pass thru $A$ or $B$ has length $<\pi$.
\end{itemize}
The polygonal lines in $\triangle$ that meet these properties will be called \emph{trajectories}.

\begin{thm}{Lemma}\label{lem:short-trajectories}
Let $\triangle$, $A,B,C$ and $M$ be as above.
Then, up to symmetry of $\triangle$ there are only the following trajectories of length less than $2\cdot\pi$:
\[CM,\quad AB,\quad ABAB,\quad ABAC,\quad AC,\quad ACAC,\quad \text{and}\quad ACACAC.\]

\end{thm}

\parit{Proof.}
Reflecting $\triangle$ in its sides gives the tetrahedral tiling of $\mathbb{S}^2$.
The images of $A$ are the vertices of a regular tetrahedron, the images of $B$ are their antipodes, and the images of $C$ are the midpoints of its edges.

Put $\alpha=|A-C|_{\mathbb{S}^2}$ and $\beta=|A-B|_{\mathbb{S}^2}$.
The spherical cosine rule gives
\[
\cos\alpha=\frac1{\sqrt3},
\qquad
\cos\beta=\frac13,
\qquad
2\cdot\alpha+\beta=\pi.
\]
In particular,
\[
\tfrac\pi4<\alpha<\tfrac\pi3,
\qquad
\tfrac\pi3<\beta<\tfrac\pi2.
\]

Choose a trajectory $\gamma$.
Cut $\gamma$ at all passages thru $A$ and $B$.
Each resulting arc has length $<\pi$, and after unfolding it is a minimizing geodesic in $\mathbb{S}^2$.
Thus, up to symmetry, it is either $AB$, of length $\beta$, or $ACA$, of length $2\alpha$.
Moreover, the number of $AB$-arcs is even.
Hence, writing it as $2\cdot k$ arcs of the first kind and $q$ arcs of the second kind, we have
\[
2\cdot k\cdot \beta+2\cdot q\cdot \alpha<2\cdot \pi.
\]
The inequalities above leave only the following possibilities for pair $(k,q)$:
\[(1,0),\quad (2,0),\quad(1,1),\quad(0,1),\quad(0,2),\quad\text{and}\quad(0,3),\]
which give respectively
\[
AB,\quad ABAB,\quad ABAC,\quad AC,\quad ACAC,\quad \text{and}\quad ACACAC.
\]

It remains to consider the case when $\gamma$ misses $A$ and $B$.
Then $\length\gamma<\pi$.
Stability forces $\gamma$ to pass through $C$; otherwise its spherical development can be moved slightly to shorten it.
Unfolding at $C$, the nearest distinct images of $C$ are at distance $\tfrac\pi2$, and, up to symmetry, the corresponding trajectory is $CM$.
\qeds

\section{Encoding}

Let $\tilde \gamma$ be a closed geodesic in $\Sigma$ such that $\length\gamma<2\cdot\pi$.
Note that $\tilde \gamma$ enters and exits interior of each copy of $\triangle$ thru different sides;
if not, we can make a shortcut along the side.
Therefore we can approximate $\tilde \gamma$ by a curve $\tilde \gamma'$ that crossed sides of triangle at finite number of points; moreover, if two points consequent points lie on one side, then these are two vertices of this side.

Recall that each triangle in $\Sigma$ is labeled by an element of $B_4$.
Let us write a cyclic sequence of labels $b_1,\dots,b_k$ of triangle with their interior visited by $\tilde \gamma'$;
now let us label $i$-th crossing point by $z_i=b_{i+1}\cdot b_i^{-1}$, so at the $i$-th crossing you multiply the old label by $z_i$.
Evidently, $z_1\cdots z_k=1$.
The sequence $z_1,\dots,z_k$ will be called \emph{encoding} of $\tilde\gamma$;
it is defined almost uniquely: if $i$-th edge of $\tilde\gamma$ runs along a side, say $c$, then changing labels $z_{i-1}$ and $z_i$ to $c^k\cdot z_{i-1}$ and $z_i\cdot c^{-k}$ produces another set of labels.

Let us list possible encodings for curves that project to a trajectories of length less than $2\cdot\pi$ in the spherical triangle $\triangle$ with an aditional constrain that the trajectory lifts to a closed geodesic.

\parit{$CM$.} The trajectory $CM$ can have an encoding $(x,y)$, where
\begin{itemize}
\item $x\in \langle a,b\rangle=\langle a\rangle\cdot\langle b\rangle$ at $C$, and since the lift is geodesic, $x\notin\langle a\rangle\cup \langle b\rangle$.
\item $y\in \langle c\rangle$ at $M$, and since the lift is geodesic, $y\ne 1$.
\end{itemize}
In other words $x=a^m\cdot b^n$ and $y=c^k$ for some $m\ne 0$, $n\ne 0$, and $k\ne 0$.
Since $a^m\cdot b^n\cdot c^k=1$ if and only if $m=n=k=0$, this trajectory does not have geodesic lifts.

\parit{$AB$.} The trajectory $AB$ can have an encoding $(x,y)$, where
\begin{itemize}
\item $x\in \langle b,c\rangle$ at $A$, but since the lift is geodesic, $x\notin \langle c\rangle\cdot\langle b\rangle\cdot\langle c\rangle$;
\item $y\in \langle a,c\rangle$ at $B$, but since the lift is geodesic, $y\notin \langle c\rangle\cdot\langle a\rangle\cdot\langle c\rangle$.
\end{itemize}
To show that this does not have geodesic lifts, we need to check that $x\cdot y\ne1\in B_4$.
If $x\cdot y=1$, then $x\cdot y\cdot x\cdot y=1$, but the latter is not possible from the case $ABAB$.

\parit{$ABAB$.} The trajectory $ABAB$ can have an encoding $(x_1,y_1,x_2,y_2)$, where
\begin{itemize}
\item $x_1,x_2\in \langle b,c\rangle$ at $A$, but since the lift is geodesic, $x_1,x_2\notin \langle c\rangle\cdot\langle b\rangle\cdot\langle c\rangle$;
\item $y_1,y_2\in \langle a,c\rangle$ at $B$, but since the lift is geodesic, $y_1,y_2\notin \langle c\rangle\cdot\langle a\rangle\cdot\langle c\rangle$.
\end{itemize}
Again, we need to check that $x_1\cdot y_1\cdot x_2\cdot y_2\ne1\in B_4$, which is proved in
Corollary \ref{cor:ABAB} of Agol's lemma.

\parit{$ABAC$.} The trajectory $ABAC$ can have an encoding $(x_1,y,x_2,z)$, where
\begin{itemize}
\item $x_1,x_2\in \langle b,c\rangle$ at $A$, but since the lift is geodesic, $x_1\notin \langle b\rangle\cdot\langle c\rangle$ and
$x_2\z\notin \langle c\rangle\cdot\langle b\rangle$;
\item $y\in \langle a,c\rangle$ at $B$, but since the lift is geodesic, $y\notin \langle c\rangle\cdot\langle a\rangle\cdot\langle c\rangle$;
\item $z\in  \langle a,b\rangle=\langle a\rangle\cdot\langle b\rangle$ at $C$, but since the lift is geodesic, $z\notin \langle b\rangle$.
\end{itemize}
We need to check that $x_1\cdot y\cdot x_2\cdot z\ne 1\in B_4$, and this also follows from the first Agol's lemma.

\parit{Case $AC$:} we have elements
\begin{itemize}
 \item $x\in \langle b,c\rangle$ at $A$, but since the lift is geodesic, $x\notin\langle b\rangle\cdot\langle c\rangle\cdot\langle b\rangle$;
 \item $z\in \langle a,b\rangle=\langle a\rangle\cdot\langle b\rangle$ at $C$, but since the lift is geodesic, $z\notin \langle b\rangle$.
\end{itemize}
Since the lift is closed in the braid space, we have $x\cdot z=1\in B_4$.
So, the problem reduces the following:
\[(\langle b,c\rangle \setminus \langle b\rangle\cdot\langle c\rangle\cdot\langle b\rangle)
\cdot
(\langle a,b\rangle \setminus \langle b\rangle)\not\ni 1.\]
This case, as well as the next will follow from the next one.

\parit{Case $ACAC$:} we have elements
\begin{itemize}
 \item $x_1,x_2\in \langle b,c\rangle$ at $A$, but since the lift is geodesic, $x_1,x_2\notin\langle b\rangle\cdot\langle c\rangle\cdot\langle b\rangle$;
 \item $z_1,z_2\in \langle a,b\rangle=\langle a\rangle\cdot\langle b\rangle$ at $C$, but since the lift is geodesic, $z_1,z_2\notin \langle b\rangle$.
\end{itemize}
Since the lift is closed in the braid space, we have $x_1\cdot z_1\cdot x_2\cdot z_2 =1\in B_4$.
So, the problem reduces the following:
\[(\langle b,c\rangle \setminus \langle b\rangle\cdot\langle c\rangle\cdot\langle b\rangle)
\cdot
(\langle a,b\rangle \setminus \langle b\rangle)
\cdot
(\langle b,c\rangle \setminus \langle b\rangle\cdot\langle c\rangle\cdot\langle b\rangle)
\cdot
(\langle a,b\rangle \setminus \langle b\rangle)
\not\ni 1;\]
equivalently,
\[(\langle b,c\rangle \setminus \langle b\rangle\cdot\langle c\rangle\cdot\langle b\rangle)
\cdot
(\langle a,b\rangle \setminus \langle b\rangle)\cap(\langle a,b\rangle \setminus \langle b\rangle)\cdot (\langle b,c\rangle \setminus \langle b\rangle\cdot\langle c\rangle\cdot\langle b\rangle)=\varnothing.\]

\parit{Case $ACACAC$:} we have elements
\begin{itemize}
 \item $x_1,x_2,x_3\in \langle b,c\rangle$ at $A$, but since the lift is geodesic, $x_1,x_2,x_3\notin\langle b\rangle\cdot\langle c\rangle\cdot\langle b\rangle$;
 \item $z_1,z_2,z_3\in \langle a,b\rangle=\langle a\rangle\cdot\langle b\rangle$ at $C$, but since the lift is geodesic, $z_1,z_2,z_3\notin \langle b\rangle$.
\end{itemize}
Since the lift is closed in the braid space, we have $x_1\cdot z_1\cdot x_2\cdot z_2\cdot x_3\cdot z_3 =1\in B_4$.
So, the problem reduces the following:
\[(\langle b,c\rangle \setminus \langle b\rangle\cdot\langle c\rangle\cdot\langle b\rangle)
\cdot
(\langle a,b\rangle \setminus \langle b\rangle)
\cdot
(\langle b,c\rangle \setminus \langle b\rangle\cdot\langle c\rangle\cdot\langle b\rangle)
\cdot
(\langle a,b\rangle \setminus \langle b\rangle)
\cdot
(\langle b,c\rangle \setminus \langle b\rangle\cdot\langle c\rangle\cdot\langle b\rangle)
\cdot
(\langle a,b\rangle \setminus \langle b\rangle)
\not\ni 1.\]
This case follows from the second Agol's lemma.

\section{Agol's lemma}


\begin{thm}{Lemma}\label{lem:agol1}
Consider the braid group $B_4$ with standard generators $a$, $b$, and $c$ as above; so $a\cdot b=b\cdot a$, $a\cdot c\cdot a=c\cdot a\cdot c$, and $b\cdot c\cdot b=c\cdot b\cdot c$.
Then
\[(\langle a,c\rangle\cdot\langle b,c\rangle)\cap (\langle b,c\rangle\cdot\langle a,c\rangle)
=
(\langle c\rangle\cdot \langle a,b\rangle\cdot\langle c\rangle)\cup\langle a,c\rangle\cup \langle b,c\rangle.\]
\end{thm}

\parit{Proof.}
Let us identify $B_4$ with the mapping class group of the 4-punctured disk  $\DD^2\setminus \{p_1,p_2,p_3,p_4\}$ fixing the boundary.

Let us choose a complete hyperbolic metric on $\DD^2\setminus \{p_1,p_2,p_3,p_4\}$ so that the boundary is totally geodesic.
Choose an arc $w$ in $\DD^2\setminus \{p_1,p_2,p_3,p_4\}$ with endpoints on $\partial \DD^2$.
In the isotopy class of $w$ there is a unique geodesic representative.
Moreover the intersection number between pairs of arcs will be minimized by these representatives.
Abusing notation, let us write $s(w)$ for the arc made geodesic after applying $s\in B_4$ to $w$.

\begin{figure}[ht!]
\centering
\vskip-0mm
\includegraphics{mppics/pic-10}
\vskip-0mm
\end{figure}

Consider the arcs $x, y, z$ as in the picture; so
\begin{align*}
\stab(y)&=\langle a,c\rangle,
&
\stab(x)&=\langle b,c\rangle,
&
\stab(x\cup y)&=\langle c\rangle.
\end{align*}

Suppose
\[s\in (\langle a,c\rangle\cdot\langle b,c\rangle)\cap(\langle b,c\rangle\cdot\langle a,c\rangle);\]
in other words, $s=h_1\cdot k_1=k_2\cdot h_2$ for some $h_1,h_2\in \langle a,c\rangle$ and $k_1,k_2\in \langle b,c\rangle$.
Then  $s(x)\cap y =\emptyset$ and $s(y)\cap x = \emptyset$.
Moreover $s(x)\cap s(y)=\emptyset$ since $x\cap y=\emptyset$.
Indeed, $k_i(x)=x$ and $h_i(y)=y$, and so
\[s(x)\cap y =h_1\cdot k_1(x)\cap y =h_1(x)\cap h_1(y)=\emptyset;\]
similarly, one gets $s(y)\cap x = \emptyset$.

\begin{figure}[ht!]
\centering
\vskip-0mm
\includegraphics{mppics/pic-21}
\vskip-0mm
\end{figure}

Consider the subdisc $M\subset \DD^2$ that is the middle piese of $\DD^2$ after cutting along $x$ and $y$.
Note that $M\ni p_2, p_3$, so $M\setminus \{p_2,p_3\}$ is a twice-punctured disk.
Note that $M\setminus \{p_2,p_3\}$ may contain only one isotopy class of nonintersecting arcs with endpoints on the boundary.
Therefore, each component of $s(x)\z\cap M$ and $s(y)\cap M$ must be isotopic to one arc.
All components of $s(x)\cap M$ must have endpoints on $x\subset \partial M$,
and all components of $s(y) \cap M$ must have endpoints on $y\subset \partial M$.%
\footnote{One of endpoints $s(x)\cap M$ must be an endpoint of $x$ and one of endpoints $s(y)\cap M$ must be an endpoint of $y$, but it is not important.}


We may assume that $s(x)\ne x$ and $s(y)\ne y$; in other words, there are components of $s(x)\cap M$ and $s(y) \cap M$.
Otherwise
\[s\in \langle b,c\rangle =\stab(x)\quad\text{and}\quad s\in \langle a,c\rangle =\stab(y),\] respectively.

Then by composing $s$ with an element of $\langle c\rangle$, we can unwind this configuration;
that is, replacing $s$ with $s'=c^{-i}\cdot s$ if necessary, we acheave that $s'(x)\cap z=\emptyset$ and
$s'(y)\cap z=\emptyset$.


\begin{figure}[ht!]
\centering
\vskip-0mm
\includegraphics{mppics/pic-30}
\vskip-0mm
\end{figure}

Now let us focus on the subdiscs $L$ and $R$ in $\DD^2$ and obtained by cutting along $z$ and taking the left and right piece, respectively.
There is only one isotopy class of arc $x$ and $s(x)$ in $L\setminus \{p_1,p_2\}$.
Thus, replacing $s'$ with $s''=a^{-j}\cdot s'$, we acheave that $s''(x)=x$.
Similarly, replacing $s''$ with $s'''=b^{-k}\cdot s''$, we acheave that $s'''(y)=x$.

Since $\stab(x \cup y) =\langle c\rangle$, we have $s'''\in \langle c\rangle$.
Therefore
$s = c^i a^j b^k c^l$.
In other words,
\[s\in \langle c\rangle\cdot \langle a,b\rangle\cdot\langle c\rangle.\]
\qeds


\begin{thm}{Corollary}\label{cor:ABAB}
Suppose
\[h_1\cdot k_1\cdot h_2 \cdot k_2=1,\]
where $h_i\in \langle a,c\rangle$ and $k_i\in \langle b,c\rangle$.
Then $h_i\in\langle a\rangle\cdot \langle c\rangle\cdot \langle a\rangle$ or $k_i\in\langle b\rangle\cdot \langle c\rangle\cdot \langle b\rangle$ for some $i$.
\end{thm}

\parit{Proof.}
Let $s=h_1\cdot k_1= k_2^{-1} \cdot h_2^{-1}$.
By \ref{lem:agol1}, we need to consider three cases

Suppose  $s \in \langle c\rangle\cdot \langle a,b\rangle\cdot\langle c\rangle$;
in other words, $s = c^i\cdot a^j\cdot b^k\cdot c^l$ for some $i$, $j$, $k$, and $l$.
Then,
\[h_1^{-1} c^i a^j = k_1 c^{-l} b^{-k} \in \langle a,c\rangle\cap \langle b,c\rangle = \stab(x\cup y)=\langle c\rangle.\]
Thus $h_1 = c^i a^j c^{-m}$ and $k_1 =c^m b^k c^l$, which already implies more than we need.

If $s \in \langle b,c\rangle$, then $h_1\in\langle a,c\rangle\cap \langle b,c\rangle=\langle c\rangle$; similarly $h_2\in\langle c\rangle$ --- again, this implies more than we need.
The last case $s \in \langle a,c\rangle$ can be done the same way.
\qeds

\section{Second lemma}

\begin{thm}{Second lemma}
Consider the braid group $B_4$ with standard generators $a=\sigma_1$, $b=\sigma_3$, and $c=\sigma_2$; so $a\cdot b=b\cdot a$, $a\cdot c\cdot a=c\cdot a\cdot c$, and $b\cdot c\cdot b=c\cdot b\cdot c$.
Let $d=a^{-1}\cdot b$.
Suppose
\[h_1\cdot d^{n_1}\cdot h_2 \cdot d^{n_2}\cdot  h_3\cdot  d^{n_3}=1,\]
where $h_i\in \langle a,c\rangle$ and $n_i\ne 0$.
Then $h_i\in\langle a\rangle\cdot \langle c\rangle\cdot \langle a\rangle$ for some $i$.
\end{thm}

\parit{Proof.}
This equation can be rewritten as
\[h_1\cdot h_2\cdot h_3\cdot x^{n_1}\cdot y^{n_2}\cdot z^{n_3}=1,\]
where $x=h_3^{-1}\cdot h_2^{-1}\cdot d\cdot h_2\cdot h_3$, $y=h_3^{-1}\cdot d\cdot h_3$, and $z=d$ are elements of the free group $F_2=\ker\phi=\llangle d\rrangle$, and $\phi$ is the homomorphism $B_4\to B_3$ that identifies $a$ and $b$.

Note that $h_1\cdot h_2\cdot h_3=1$ and $x^{n_1}\cdot y^{n_2}\cdot z^{n_3}=1$.
Furthermore, $x$, $y$ and $z$ are primitive elements of $F_2$.

We can assume that the subgroups $\langle y, z\rangle$ is not cyclic.
Otherwise $h_3\in\langle a\rangle$, and the statent follows.

By the Lyndon--Schützenberger theorem, we can assume that $|n_1|=1$.

Furthermore, we can assume that $|n_1|=|n_2|=1$.
Indeed, assume $|n_2|\ge 2$ and $|n_3|\ge 2$.
Note that $x\in \langle y,z\rangle$.
Since $x$ is primitive in $F_2$, by the Kurosh subgroup theorem $x$ is primitive in $\langle y,z\rangle$.
Therefore $\ZZ_{n_2}*\ZZ_{n_2}\cong \langle y,z\rangle/\llangle y^{n_2},z^{n_3}\rrangle$ is a quotient of $\langle y,z\rangle/\llangle y^{n_2}\cdot z^{n_3}\rrangle\cong \langle y,z\rangle/\llangle x \rrangle\cong \ZZ$,
which is impossible since $\ZZ_{n_2}*\ZZ_{n_2}$ is not commutative.

To emphasise that $|n_1|=|n_2|=1$, put $\eps_1=n_1$, $\eps_2=n_1$, and $N=n_3$.
Let us rewrite the original equation as
\[h_1\cdot d^{\eps_1}\cdot h_2 = d^{-N}\cdot h_3^{-1}\cdot d^{-\eps_2}.\]
Arguing as in Agol's lemma, we get
\[h_1\cdot\langle a\rangle= a^{-N}\cdot h_3^{-1}\cdot\langle a\rangle.\]
Indeed ...


Therefore,
\[
h_1=a^{-N}\cdot h_3^{-1}\cdot a^u
\quad\text{and}\quad
x=a^{-N}\cdot y\cdot a^N.
\]
for some $u\in\mathbb Z$.
Let $e=d^c$.
On $H_1(F_2)=\mathbb Z[d]\oplus\mathbb Z[e]$, conjugation by $a$
acts by $(r,s)\mapsto(r-s,s)$.
Thus, writing $[y]=(r,s)$, the relation
$x^{\epsilon_1}\cdot y^{\epsilon_2}\cdot d^N=1$ gives
\[
\epsilon_1\cdot (r-Ns,s)+\epsilon_2\cdot (r,s)+N\cdot (1,0)=0.
\]

If $\epsilon_1=\epsilon_2$, then $s=0$.
Hence $h_3\in\langle a\rangle\cdot Z(B_3)$.
Suppose $\zeta=(a\cdot c)^3$; it generates $Z(B_3)$.
Note that $h_3=a^r\cdot\zeta^k$ and $h_1=a^q\cdot \zeta^{-k}$.
Therefore
\[
h_2=h_1^{-1}\cdot h_3^{-1}\in\langle a\rangle,
\]
which finishes the proof.

If $\epsilon_1=-\epsilon_2$, then $s=\epsilon_1$.
Under the standard homomorphism
$B_3\to \SL_2(\mathbb Z)$ this implies
\[
h_3=h\cdot \zeta^{2\cdot k}
\]
for some $h=a^r\cdot  c^{\epsilon_1}\cdot a^t$.
Consequently
\[
h_2
=a^{-u}\cdot h_3\cdot a^N\cdot h_3^{-1}
=a^{-u}\cdot h\cdot a^N\cdot h^{-1}
=a^{r-u-\epsilon_1}\cdot c^N\cdot a^{\epsilon_1-r},
\]
where we used
\[
c^\epsilon\cdot  a^N\cdot  c^{-\epsilon}
=a^{-\epsilon}\cdot c^N\cdot a^\epsilon,
\qquad \epsilon=\pm1.
\]
Thus, $h_2\in\langle a\rangle\cdot \langle c\rangle\cdot \langle a\rangle$, and the proof is finished.
\qeds

{\sloppy
\def\emph{\textit}
\printbibliography[heading=bibintoc]
\fussy
}

\end{document}

\section{Ramifications}

Consider a finite orientation-preserving action $\Gamma\acts\RR^m$.
Denote by $\mathcal{L}_\Gamma$ the arrangement of all hyperlines which are fixed by at least one non-identical element in $\Gamma$. Define the ramification of $\Gamma\acts\RR^m$
(briefly $\Ram_\Gamma$)
as the universal cover of $\RR^m$ branching infinitely  along each hyperline in $\mathcal{L}_\Gamma$.

More precisely, if $\tilde W_\Gamma$ denotes the universal cover of
$$W_\Gamma=\RR^m\backslash
\left(\bigcup_{\ell\in\mathcal{L}_\Gamma}\ell\right)$$
equipped with the length metric induced from $\RR^m$
then $\Ram_\Gamma$ is the metric completion of $\tilde W_\Gamma$.

We want to show the following that $\Ram_\Gamma$ is a $\CAT[0]$ for any properly discontinuous isometric orientation-preserving action $\Gamma\acts\RR^m$.

We thaut about this question for many years and I do not see why it should be true, except we cannot find a counterexample.
Let me state what we know.

Assume $\Gamma\acts\RR^m$ is a counterexample; we can assume that $m$ is minimal.
In this case the action cannot split.

Note that $\Ram_\Gamma$ is a cone.
Consider the unit sphere $S$ in $\Ram_\Gamma$;
observe that $S$ is not $\CAT[1]$.

Since $m$ is minimal, one can see that $S$ is locally $\CAT[1]$.
It follows that there is a closed geodesic $s$ in $S$ of length $<2\cdot\pi$.
Note that projection of $s$ in the unit sphere in $\RR^m$ must lie in an open hemisphere...
